\chapter{The heat trace, the arithmetic and the limits of finite information}\label{chap:heat}\chaptermark{The heat trace and finite information}

Chapter~\ref{chap:hcm} worked at the top of the Bondesson hierarchy, with hyperbolic complete monotonicity, which for a Laplace transform is the GGC property; here we climb the hierarchy from the bottom. Every rung is a sign condition on the heat trace $W(t)$ of Section~\ref{ch:clock}, the pairing of the Thorin functional with $e^{-tz}$ (not the Mellin variable $W$ of Section~\ref{ch:boschB}). The lower rungs are unconditional theorems, proved with a finite verification of the zeros; the top rung is RH. Section~\ref{ch:heat} develops this hierarchy of positivity, and Section~\ref{ch:finite} shows what such finite information cannot decide. Section~\ref{ch:arith} computes the Thorin functional from the primes, and Section~\ref{ch:spacing} relates the errors of the P\'olya route on a window of finite height to the spacings of the zeros, which it compares with random matrix theory.

Taken together, Chapters~\ref{chap:measure} to \ref{chap:heat} show that on every window where the zeros are verified the computed Thorin measure converges to the RH measure at explicit rates, while above the computed range every finite computation of the kinds in this book remains consistent with a failure of the Thorin condition (Section~\ref{ch:finite}). Three of their threads return in Chapters~\ref{chap:basepoint} and \ref{chap:almost}: the count at resolution one half, which is the basepoint one (Section~\ref{ch:basepoint}), the self decomposability of the clock (Section~\ref{ch:sato}), and the zero density, whose classical estimates bound the negative Thorin mass below the basepoint one (Section~\ref{ch:almost}).
